\section{Numerical error: complete proofs}
\label{sec:annex-p12-numerical-error}

\begin{annexitem}{annex:p12.numerical.nearest}{thm:p12-numerical-nearest}{Spacing, nearest rounding and the range boundary}
\begin{proof}
In \([\beta^e,\beta^{e+1}]\), consecutive format values are separated by
\(h=\beta^{e-p+1}\), including the upper endpoint. Every real value
between them lies within \(h/2\) of an endpoint. At the lower radix-interval
boundary the gap on the other side is smaller, so the same upper bound
still applies. Divide by \(|z|\geq\beta^e\) to obtain
\(h/(2|z|)\leq\beta^{1-p}/2=u\). Symmetry proves the negative case.
As \(|z|\) ranges from \(\beta^e\) to \(\beta^{e+1}\), the ratio
\(h/(2|z|)\) ranges from \(u\) down toward \(u/\beta\).
In a normal-only finite format, every sufficiently small positive value
rounds to zero. Its relative error is one, so the preceding small uniform
bound cannot extend to that region. Zero is handled separately.
\end{proof}
\end{annexitem}

\begin{annexitem}{annex:p12.numerical.no-guard}{thm:p12-numerical-no-guard}{The complete no-guard example}
\begin{proof}
Take \(x=1\) and \(y=1-\beta^{-p}\), whose normalized significand
has \(p\) digits all equal to \(\beta-1\). Aligned at the exponent of
\(x\), its final digit is discarded by the no-guard procedure, leaving
\(\bar y=1-\beta^{1-p}\). Thus the exact difference is
\(\beta^{-p}\), whereas the computed difference is
\(\beta^{1-p}\), already representable. Their relative discrepancy is
\((\beta^{1-p}-\beta^{-p})/\beta^{-p}=\beta-1\).
For the decimal example, alignment replaces \(9.93\) by \(9.9\),
so \(10.1-9.9=0.2\), instead of \(0.17\).
\end{proof}
\end{annexitem}

\begin{annexitem}{annex:p12.numerical.guard-bound}{thm:p12-numerical-guard-bound}{All exponent and carry cases for a guard digit}
\begin{proof}
For subtraction interchange the positive operands if necessary and scale
them by an exact power of \(\beta\), so \(1\leq x<\beta\) and
\(0<y<x\). Such scaling preserves relative error. If their exponents
agree, the difference is an integer multiple of \(\beta^{1-p}\)
smaller than \(\beta\), hence is exactly representable. If the exponents
differ by one, all digits of the smaller operand fit in the \(p+1\)
aligned positions. The operation is then exact subtraction followed by
nearest rounding, with relative error at most \(u\).

For exponent difference at least two, \(y<1/\beta\). Let \(\bar y\)
be the truncated operand and \(r=y-\bar y\), so
\(0\leq r<\beta^{-p}=2u/\beta\). Put
\(d=x-y\) and \(\bar d=x-\bar y=d+r\).
If \(d\geq1\), then \(1\leq\bar d<\beta\), and the last rounding
has absolute error at most \(u\). Therefore its total relative error
is less than \(u+2u/\beta\).

If \(d<1\), then \(\bar d<1+\beta^{-p}\). Both \(x\) and
\(\bar y\) are multiples of \(\beta^{-p}\), so
\(\bar d\leq1\). Every such nonnegative lattice value at most one
is exactly representable; the final rounding adds no error. Since
\(d=x-y>1-1/\beta\), the relative error is less than
\[
 \frac{\beta^{-p}}{1-1/\beta}=\frac{2u}{\beta-1}
 \leq (1+2/\beta)u.
\]
The last inequality is equality at \(\beta=2\); after multiplying
by \(\beta(\beta-1)\), it reduces to
\(0\leq(\beta-2)(\beta+1)\). This includes the case where truncation
moves a difference across a radix-interval boundary.

For same-sign addition normalize \(x\geq y\geq0\) in the same way,
and first use no guard digits. Set \(U=\beta^{1-p}=2u\).
Truncation loses \(r<U\); the aligned sum \(s=x+\bar y\) is a
multiple of \(U\). If \(s<\beta\), it is exactly representable,
so relative error is less than \(U\), since \(x+y\geq1\).
If \(s\geq\beta\), its final rounding has absolute error at most
\(\beta U/2\), and \(x+y\geq\beta\). Thus the relative error is
less than \(U/\beta+U/2=(1+2/\beta)u\leq2u\).
With \(g\geq1\) guard digits the discarded tail is less than
\(U/\beta^g\). Without a carry the final rounding error is at most
\(U/2=u\), giving relative error less than
\(2u/\beta^g+u\leq2u\). With a carry, divide the combined
absolute bound \(U/\beta^g+\beta U/2\) by a sum at least
\(\beta\), giving less than \(2u/\beta^{g+1}+u\leq2u\).
For subtraction with more than one guard digit, the equal-exponent
and adjacent-exponent cases still round an exact difference. In the
remaining case \(y<1/\beta\), the discarded tail is less than
\(U/\beta^g\). If the exact difference is at least one, its
relative error is bounded by \(u+2u/\beta^g\leq2u\).
Below one the rounded difference stays at most one, and its final
absolute rounding error is at most \(u/\beta\). For \(g\geq2\)
the relative error is therefore less than
\[
 \frac{u/\beta+2u/\beta^g}{1-1/\beta}
 \leq\frac{\beta+2}{\beta(\beta-1)}u\leq2u.
\]
This proves the at-least-one-guard convention used by later formulas. Signed cases
reduce to these two magnitude operations. Equal magnitudes subtract
exactly to zero, and adding zero is exact.
\end{proof}
\end{annexitem}

\begin{annexitem}{annex:p12.numerical.close-subtraction}{thm:p12-numerical-close-subtraction}{The factor-two exactness argument}
\begin{proof}
Take \(x\geq y>0\); the other ordering changes only the sign. If the
exponents agree, the difference is a multiple of their shared last-place
unit, with no more than \(p\) significant digits. Otherwise
\(x\leq2y\) and \(\beta\geq2\) imply that the exponents differ by
at most one. Normalize \(x\) to exponent zero; then
\(y=0.d_1\cdots d_p\) in the aligned notation. A guard digit retains
this entire value. Moreover \(0\leq x-y\leq y<1\), and the exact
difference is a multiple of \(\beta^{-p}\). It therefore has at most
\(p\) significant digits and the final rounding changes nothing.
The finite-format transfer requires that this exact difference not be
lost at the exponent boundary.
\end{proof}
\end{annexitem}

\begin{annexitem}{annex:p12.numerical.product-errors}{thm:p12-numerical-product-errors}{Exact product bounds without dropping terms}
\begin{proof}
Expand the product minus one as a sum over nonempty subsets of
\(\{1,\ldots,n\}\). The triangle inequality bounds every term
\(\prod_{i\in S}\delta_i\) by \(\prod_{i\in S}a_i\).
Summing these terms gives exactly \(\prod_i(1+a_i)-1\).
For the difference-of-squares evaluation the two signed operations have
errors bounded by \(2u\), and the multiplication by \(u\).
Multiplying their three relative factors and expanding gives the
displayed polynomial. If an exact sum or difference is zero, the
corresponding zero computation is exact; the theorem only asserts a
relative bound when the target product is nonzero.
\end{proof}
\end{annexitem}

\begin{annexitem}{annex:p12.numerical.cancellation}{thm:p12-numerical-cancellation}{Operand error survives exact subtraction}
\begin{proof}
The difference between \(\widehat A-\widehat B\) and \(A-B\) is
\(e_A-e_B\), whose magnitude is at most \(E_A+E_B\).
Add any subtraction error by the triangle inequality. Division by
\(|A-B|\) gives the relative statement when this denominator is nonzero.
Choosing \(A-B\) arbitrarily small while the operand-error budget
stays fixed proves that no small uniform relative bound follows.
\end{proof}
\end{annexitem}

\begin{annexitem}{annex:p12.numerical.quadratic}{thm:p12-numerical-quadratic}{Conjugate root identities and their denominator conditions}
\begin{proof}
Since \(b>0\) and \(\sqrt d\geq0\), the denominator
\(b+\sqrt d\) is nonzero. Multiply the direct numerator by its
conjugate:
\[
 \frac{-b+\sqrt d}{2a}
 =\frac{d-b^2}{2a(b+\sqrt d)}
 =-\frac{2c}{b+\sqrt d}.
\]
This also covers \(c=0\). When \(b<0\), use
\((-b-\sqrt d)/(2a)=-2c/(b-\sqrt d)\), whose conjugate
denominator is then negative and nonzero. These are real identities;
they make no claim about an already perturbed discriminant or coefficients.
\end{proof}
\end{annexitem}

\begin{annexitem}{annex:p12.numerical.triangle-radicand}{thm:p12-numerical-triangle-radicand}{All four factors, multiplication errors and degeneracy}
\begin{proof}
The ordering and triangle inequality give \(b\leq a\leq2b\), so
\(a\ominus b=a-b\) exactly by the close-subtraction theorem.
The first factor \(a\oplus(b\oplus c)\) differs multiplicatively
from \(a+b+c\) by two factors with errors at most \(2u\): the inner
sum's perturbation is attenuated by \((b+c)/(a+b+c)\leq1\), and
the outer sum contributes another factor. The second and third factors
each contribute one error of magnitude at most \(2u\), because their
inputs \(c,a-b\) are exact. The fourth factor
\(a\oplus(b\ominus c)\) contributes two such factors; the inner
perturbation is attenuated by \((b-c)/(a+b-c)\leq1\).
Zero inner differences introduce no error. All four exact factors are
positive in a nondegenerate triangle. Three multiplications contribute
three factors with error at most \(u\). Therefore
\[
 (1-2u)^6(1-u)^3\leq Q/q\leq(1+2u)^6(1+u)^3.
\]
For \(0\leq u\leq1/200\), the polynomial
\(\{(1+2u)^6(1+u)^3-1\}/u\), extended continuously at zero,
has nonnegative coefficients and is increasing. At \(u=1/200\),
\[
 (1+2u)^6(1+u)^3
 =\frac{8620181596490631201}{8000000000000000000}
 <\frac{27}{25}=1+16u.
\]
Thus the upper relative error is less than \(16u\) throughout the
stated interval. Repeated use of \((1-r)(1-s)\geq1-r-s\) for
\(r,s\in[0,1]\) bounds the lower product by \(1-15u\).
For \(a=b+c\), the exact inner difference is \(a-b=c\), so
the second factor is zero exactly, and the radicand is zero.
\end{proof}
\end{annexitem}

\begin{annexitem}{annex:p12.numerical.triangle-area}{thm:p12-numerical-triangle-area}{Area identity and propagation through square root and division}
\begin{proof}
Let \(s=(a+b+c)/2\). The four real factors are
\(2s,2(s-a),2(s-b),2(s-c)\), so
\(q=16s(s-a)(s-b)(s-c)\). To verify the area identity without
assuming Heron's formula, use the included angle between sides \(b,c\).
Its cosine is \((b^2+c^2-a^2)/(2bc)\), from the squared distance
formula; area is \(bc\sin\theta/2\). Squaring gives
\(16A^2=4b^2c^2-(b^2+c^2-a^2)^2=q\).
The nonnegative square root therefore yields \(A=\sqrt q/4\).

Write \(Q=q(1+\delta)\), with \(|\delta|<16u\leq0.08\).
Rationalizing the square-root difference gives
\[
 |\sqrt{1+\delta}-1|
 =\frac{|\delta|}{\sqrt{1+\delta}+1}
 \leq0.52|\delta|<8.5u,
\]
since \(\sqrt{0.92}+1>1/0.52\), verified by squaring positive
quantities. The half-ulp square-root assumption and nearest rounding
of division by four add two factors bounded by \(u\). The resulting upper factor
is \((1+8.5u)(1+u)^2\), whose error divided by \(u\) is increasing.
At \(u=1/200\) this factor is \(16847217/16000000<211/200\),
which is \(1+11u\). The lower factor is at least \(1-10.5u\)
by the same product inequality used above. Binary division by four is
an exponent shift and can omit its rounding factor. A degenerate
triangle has \(q=Q=0\), so both areas are zero.
\end{proof}
\end{annexitem}

\begin{annexitem}{annex:p12.numerical.mu}{thm:p12-numerical-mu}{Integral bounds for the logarithmic ratio}
\begin{proof}
Integrating \((1+xt)^{-1}\) for \(x>0\) proves the representation;
at zero the integral is one. Differentiation under this continuous
integrand on a compact interval gives
\(\mu'(x)=-\int_0^1t(1+xt)^{-2}\,dt\).
It follows that \(-1/2\leq\mu'(x)\leq0\), and that
\(\mu(x)\geq1/(1+x)\), hence \(\mu(x)\geq1/2\) on \([0,1]\).
For the logarithmic derivative put \(w(t)=(1+xt)^{-1}\).
This function is nonincreasing and lies in \((0,1]\). Thus
\[
 -\mu'(x)=\int_0^1t w(t)^2\,dt
 \leq\int_0^1t w(t)\,dt\leq\tfrac12\int_0^1w(t)\,dt.
\]
For the last inequality pair \(t\) with \(1-t\): the integral of
\((t-1/2)w(t)\) is
\(\int_0^{1/2}(t-1/2)(w(t)-w(1-t))\,dt\leq0\).
Division by the positive \(\mu(x)\) proves the claim. These integral
arguments also establish the endpoint derivatives by continuity.
\end{proof}
\end{annexitem}

\begin{annexitem}{annex:p12.numerical.log-increment}{thm:p12-numerical-log-increment}{Guard-digit addition and the compensated logarithm}
\begin{proof}
Normalize the addition \(1+x\) at exponent zero. Its one-guard
truncation discards less than \(\beta^{-p}=2u/\beta\), and the
final rounding, still in the radix interval starting at one, adds absolute
error at most \(u\). Consequently
\(z=1+\widehat x\), where \(\widehat x\geq0\) and
\(|\widehat x-x|\leq2u\). Also \(z<2\): before rounding the sum
is less than \(1.75\), and the upper rounding displacement is less
than \(0.1\). The close-subtraction theorem therefore gives
\(t=z\ominus1=z-1=\widehat x\) exactly.

If \(z=1\), then \(0<x\leq2u\). The elementary inequality
\(\ln(1+x)\geq2x/(2+x)\) follows because the difference is zero
at zero and its derivative is
\(x^2/((1+x)(2+x)^2)\geq0\). Thus returning \(x\) has relative
error at most \(x/2\leq u\).

Otherwise \(\widehat x>0\). The logarithmic-ratio lemma and the
mean value theorem give
\[
 \exp(-u)\leq\mu(\widehat x)/\mu(x)\leq\exp(u).
\]
The half-ulp library logarithm, nearest-rounded division and final
multiplication contribute three factors, each between \(1-u\) and \(1+u\).
Hence
\[
 \exp(-u)(1-u)^3\leq L/\ln(1+x)\leq\exp(u)(1+u)^3.
\]
For \(0\leq u<0.1\), the power series yields
\(\exp(u)\leq1/(1-u)\). Further,
\((1+u)^3/(1-u)\leq1+5u\), since after multiplication by
\(1-u\) the difference is \(u(1-8u-u^2)\geq0\).
For the lower bound, \(\exp(-u)\geq1-u\), so the lower factor is
at least \((1-u)^4\geq1-4u\). This proves the five-unit bound
without silently upgrading one-guard addition to exact rounding.
The branch at zero returns zero without a relative-error claim.
\end{proof}
\end{annexitem}

\begin{annexitem}{annex:p12.numerical.no-drift}{thm:p12-numerical-no-drift}{Symmetric rounding cells and the exact asymmetric-boundary case}
\begin{proof}
Every nonzero format value has immediate neighbors. Away from a signed
radix power their two gaps are the same power of \(\beta\), say
\(D\). Its nearest-rounding cell has endpoints \(z\pm D/2\),
both included when its retained integer significand is even and both
excluded when odd. This follows from the equally spaced digit grid
and alternating parity of neighboring coefficients. At a positive
boundary \(z=\beta^k\), put
\[
 q_0=\beta^{k-p+1},\qquad a_0=q_0/\beta.
\]
The left gap is \(a_0\), the right gap is \(q_0\), and the
rounding cell is \([z-a_0/2,z+q_0/2]\). Both endpoints belong:
the boundary coefficient \(\beta^{p-1}\) is even, and the
neighboring coefficients \(\beta^p-1\) and \(\beta^{p-1}+1\)
are odd. Reflection gives the negative cells. Rounding is
sign-symmetric for the same reason. In the unrestricted-exponent
model only zero rounds to zero: a nonzero real number has a
nonzero format value of its sign closer than zero.

Fix any signed format values \(x,y\) and set
\[
 z=R(x-y),\quad \rho=z-(x-y),\quad s=z+y=x+\rho,
 \quad x'=R(s),\quad\eta=x'-s.
\]
Because the format value \(x\) is itself a candidate when rounding
\(s\), nearest rounding gives \(|\eta|\leq|\rho|\).
The exact next subtraction is \(x'-y=z+\eta\).
If \(x'=x\), the two-operation map already fixes its original
input. If \(z=0\), then \(x=y\), which is such an immediate case.

Suppose \(z\) has equal adjacent gaps \(D\). The first rounding
gives \(|\rho|\leq D/2\), strictly if \(z\) is odd. Thus
\(z+\eta\) lies in the same rounding cell, including the correct
endpoint convention. The next subtraction returns \(z\), and
the next addition returns \(R(z+y)=x'\).

For an unequal-gap value, sign symmetry permits \(z=\beta^k>0\),
without restricting either original operand's sign. If \(\rho\geq0\),
its first rounding cell gives \(\rho\leq a_0/2\). Then
\(|\eta|\leq a_0/2\), and again \(z+\eta\) rounds to \(z\).
Otherwise write \(\rho=-r\), so
\[
 0<r\leq q_0/2,\qquad s=x-r,\qquad |\eta|\leq r.
\]
If \(\eta\geq-a_0/2\), the cell still contains \(z+\eta\).
Only \(\eta<-a_0/2\) remains.

In this exceptional case let \(t_-=x'=R(s)<s\) and let
\(t_+\) be the next format value above \(t_-\). Then
\(t_-<s<t_+\): an intervening or equal format value would be
closer than \(t_-\). Neither zero nor its accumulation causes a
gap here, since \(R(s)=0\) would imply \(s=0\) and \(\eta=0\).
Put \(D=t_+-t_-\) and \(d=s-t_-=-\eta\).
Nearest rounding implies \(d\leq D-d\). Since \(d>a_0/2\),
we have \(D>a_0\). Consecutive-format gaps are radix powers,
so \(D\geq\beta a_0=q_0\). Moreover \(x=s+r>s\) is a format
value, so \(x\geq t_+\). Hence
\[
 \frac{q_0}{2}\geq r=x-s\geq t_+-s
 \geq\frac D2\geq\frac{q_0}{2}.
\]
Every inequality is equality: \(D=q_0\), \(r=d=q_0/2\),
\(x=t_+\), and \(\eta=-q_0/2\). The exact next difference is
\[
 v=z-q_0/2=(\beta^p-\beta/2)\beta^{k-p}.
\]
Its coefficient is an integer in
\([\beta^{p-1},\beta^p)\), since \(\beta\) is even and
\(p\geq2\); therefore \(v\) is exactly representable.
The next subtraction returns \(v\), and its following addition
is \(v+y=s-q_0/2=t_-=x'\), also exact. Thus this last case fixes
\(x'\) as well, even though its rounded intermediate difference
can change. The two-operation map is idempotent in every case.
Induction gives \(x_j=x_1\) for all \(j\geq1\), and if
\(x_1=x\), it is constant from the beginning.
\end{proof}
\end{annexitem}

\begin{annexitem}{annex:p12.numerical.split-rounding}{thm:p12-numerical-split-rounding}{Integer lattices, midpoint parity and the additional carry case}
\begin{proof}
Zero is immediate, and sign symmetry reduces to \(x>0\). Scale to
an integer \(X\) with \(P/\beta\leq X<P\), where
\(P=\beta^p\), and write \(B=\beta^k\). All rounding operations
commute with this exact radix scaling. Set
\(T=R((B+1)X)=BV\). The integer \(V\) exists because the spacing
at the product is \(B\) when \((B+1)X<BP\), and \(\beta B\)
otherwise. If the product is below \(BP\), its rounding error is
at most \(B/2\); at or above \(BP\), it is at most \(\beta B/2\).
The product is less than \(\beta BP\), since \(B+1<\beta B\).

Put \(Z=T-X\). First \(Z\geq BP/\beta\). In the no-carry
product case this follows from \(Z\geq BX-B/2\) when
\(X>P/\beta\); at \(X=P/\beta\), the product and subtraction
are exact, because \(B\) divides \(X\). In the carry case,
\[
 X\geq BP/(B+1),\qquad Z\geq B(X-\beta/2).
\]
For \(\beta\geq4\), using \(P\geq\beta B\) gives
\[
 \frac{BP}{B+1}-\frac P\beta
 \geq\frac{B((\beta-1)B-1)}{B+1}\geq B\geq\beta/2.
\]
For \(\beta=2,B\geq4\), the same lower bound is
\(B(B-1)/(B+1)\geq1\). If \(\beta=B=2\) and \(p\geq3\),
the original difference is \(P/6\geq4/3\). The only remaining
case is \(\beta=2,p=2,k=1\), where carry forces \(X=3\),
and \(X-\beta/2=2=P/\beta\). This proves the lower bound in
every permitted parameter case.

Suppose first that \(Z<BP\). It then lies in a radix interval whose
spacing is \(B\), so
\[
 R(Z)=B\,\operatorname{nint}(V-X/B),
\]
where \(\operatorname{nint}\) is nearest-integer rounding to even.
If \(X/B\) is not a half-integer, translation and reflection give
\(\operatorname{nint}(V-X/B)=V-\operatorname{nint}(X/B)\).
If it is a half-integer, \(V\) is even: without a product carry,
\(V\) is itself the rounded integer of the half-integer
\(X+X/B\); with a carry, \(V\) is a multiple of the even
\(\beta\). Translation by an even integer preserves midpoint
parity, so the same identity holds. Consequently
\(T-R(Z)=B\operatorname{nint}(X/B)\).

It remains to handle \(Z\geq BP\), which can occur after a
product carry. Since \(X\leq P-1\),
\[
 BP\leq Z\leq BX+\beta B/2<BP+\beta B/2.
\]
Here the spacing is \(\beta B\), so \(R(Z)=BP\), strictly before
the next midpoint. Let \(C=P/B\) and \(d=P-X\).
The inequality \(Z\geq BP\) and the product-error bound imply
\(1\leq d\leq\beta/2\). Hence
\(C-1/2\leq X/B<C\), because \(B\geq\beta\).
The integer \(C=\beta^{p-k}\) is even, so
\(\operatorname{nint}(X/B)=C\), including the midpoint.
Furthermore \(V\) is a multiple of \(\beta\),
\(V-P\geq X/B>C-1\), and \(V\leq P+C\), because the
unrounded scaled product \(X+X/B\) is smaller than the representable
grid point \(P+C\). Thus \(V=P+C\). It follows that
\(T-R(Z)=B(V-P)=BC=B\operatorname{nint}(X/B)\).

In both cases this final difference has at most \(p-k\) significant
digits after normalization, including any carry, so its last rounding
is exact. Undo the scaling and sign change to obtain \(H_k(x)\).
\end{proof}
\end{annexitem}

\begin{annexitem}{annex:p12.numerical.split}{thm:p12-numerical-split}{Short residuals in the even-precision and odd-binary cases}
\begin{proof}
Use the preceding integer normalization. Write
\(X=qB+r\), \(0\leq r<B\). The computed high part is either
\(qB\) or \((q+1)B\), selected by nearest-integer rounding, and
therefore has at most \(p-k\) significant digits. The real residual
is correspondingly \(r\) or \(r-B\); its magnitude is at most
\(B/2\). At equality it is a signed value \(\beta^k/2\), which
has one significant radix digit because \(\beta\) is even.
For even \(p\), \(k=p-k\), and every smaller integer residual has
at most \(k\) radix digits. For odd binary precision,
\(k=(p+1)/2\), and an integer of magnitude strictly smaller than
\(2^{k-1}\) has at most \(k-1=(p-1)/2\) binary digits. Thus
the residual is exactly representable in all stated cases, and
\(R(x-h)=x-h\). This proves the split and its exact recombination.
Products of two \(r_0=\lfloor p/2\rfloor\)-digit significands
have at most \(2r_0\leq p\) significant digits. They are therefore
computed exactly when in range. Distributivity in the reals gives
the four-term expansion; it does not specify a lossless summation
algorithm for those terms.
\end{proof}
\end{annexitem}

\begin{annexitem}{annex:p12.numerical.round-trip}{thm:p12-numerical-round-trip}{Midpoint separation in every declared even radix}
\begin{proof}
If \(m=0\), the result is zero. Sign symmetry handles negative
\(m\). If \(i=j\), division is a radix scaling followed by division
by two. Since \(m\) has at most \(p-1\) radix digits and the base
is even, \(m/2\) has at most \(p\) digits and is exactly
representable. Multiplying back is exact. Otherwise, scale \(n\) by a radix power
so that
\(n'=\beta^{p-1}+\beta^k\), \(0\leq k\leq p-2\). The lower bound
on \(k\) follows from exact representability of \(n\) at precision
\(p\). Independently scale the positive integer \(m\) upward
by a radix power to obtain \(m'\) with
\(1/\beta\leq q=m'/n'<1\). Such a scaling is possible since
\(m<\beta^{p-1}\leq n'\). It makes \(m'\) an integer. If
\(m'\geq\beta^{p-1}\), at least one upward shift was used, so
\(m'\) is divisible by the even \(\beta\). When \(q=1/\beta\), division is exact and there is
nothing further to prove.

For \(1/\beta<q<1\), the grid spacing is \(h=\beta^{-p}\).
Every midpoint has the form \(Q=N/(2\beta^p)\) with \(N\) odd.
Then
\[
 |Q-q|=
 \frac{|N(\beta^{p-1-k}+1)-2m'\beta^{p-k}|}{2n'\beta^{p-k}}
 \geq\frac{\beta^{k-p}}{2n'},
\]
because the numerator is an odd integer and is nonzero. The distance
from \(q\) to its nearest grid point is at most half a grid step
minus its distance to the midpoint on that side. Therefore
\[
 n'|R(q)-q|
 \leq\frac{n'}{2\beta^p}-\frac{\beta^k}{2\beta^p}
 =\frac1{2\beta}.
\]
Here \(\beta^{p-2}<m'<\beta^p\). Between the lower two radix powers the
spacing is \(1/\beta\), and \(m'\), being an integer, has even
grid coefficient \(\beta m'\). At the upper radix interval its spacing is one,
and \(m'\) is even; at the shared boundary the smaller adjacent gap
is \(1/\beta\). Thus every value within \(1/(2\beta)\) of \(m'\)
rounds to \(m'\), including a possible tie on the closer side.
Apply this to \(n'R(q)\). Nearest rounding commutes with each
radix-power scaling, so undoing the independent scalings of
\(m\) and \(n\) gives the required original identity.
\end{proof}
\end{annexitem}

\begin{annexitem}{annex:p12.numerical.root-sensitivity}{thm:p12-numerical-root-sensitivity}{Square-root perturbations at and away from zero}
\begin{proof}
For \(s\geq t\geq0\),
\(s-t=(\sqrt s-\sqrt t)(\sqrt s+\sqrt t)
\geq(\sqrt s-\sqrt t)^2\).
Take square roots and interchange the two inputs if necessary to
obtain the first bound. If both inputs are at least \(d_0\),
rationalization gives denominator
\(\sqrt{\widehat d}+\sqrt d\geq2\sqrt{d_0}\), proving the
second. In either quadratic-root branch, a change only in the square
root is divided by \(2a\); take absolute values. These statements
do not define a real square root for a negative input.
\end{proof}
\end{annexitem}

\begin{annexitem}{annex:p12.numerical.certification}{thm:p12-numerical-certification}{Interval certificates and non-certification}
\begin{proof}
The assumed absolute bound is precisely
\(v\in[\widehat v-E,\widehat v+E]\). An interval strictly above
or below zero proves the corresponding sign. If it contains zero,
zero itself is compatible with the bound, so neither strict sign
follows from that bound alone. Applying the two enclosures and the
strict endpoint inequality gives the ordering result.
Finally \(|v|\geq|\widehat v|-E>0\) by the reverse triangle
inequality; division gives the relative bound. Strictness of the
endpoint conditions accounts for the boundary-equality cases.
\end{proof}
\end{annexitem}

\begin{annexitem}{annex:p12.numerical.data-budget}{thm:p12-numerical-data-budget}{The complete coefficient-uncertainty budget}
\begin{proof}
Write \(b_*=b+\Delta_b\), \(a_*=a+\Delta_a\) and
\(c_*=c+\Delta_c\). Expansion gives
\[
 b_*^2-b^2=2b\Delta_b+\Delta_b^2,\qquad
 a_*c_*-ac=a\Delta_c+c\Delta_a+\Delta_a\Delta_c.
\]
Bound each term with the stated nonnegative radii and apply the
triangle inequality, including the factor four. A further triangle
inequality adds the independent arithmetic-error budget. The two
terms concern different differences and neither bounds the other.
\end{proof}
\end{annexitem}
